% Frühe Baumaxiome; nach b27-tree-early-foundations einfügen.
\section{Das Axiomsystem der Baumstrukturen}

Die soeben bewiesenen Eigenschaften motivieren den folgenden
Strukturbegriff. Darin besitzt jedes Blatt eine Beschriftung in
\(A\), und jeder innere Knoten hat genau zwei geordnete Kinder. Die
Reihenfolge der Argumente von \(k\) entspricht der Links-rechts-Ordnung
des planaren Baumes. Eine Vertauschung wird durch die Axiome nicht
identifiziert.

Die Symbole beschreiben ganze Bäume und ihre Konstruktion:
\begin{center}
\begin{tabularx}{\linewidth}{@{}p{0.20\linewidth}X@{}}
\(A;\ a,b\) & \(A\) ist das Alphabet der Blattbeschriftungen;
  \(a\) und \(b\) bezeichnen Elemente von \(A\).\\[3pt]
\(T;\ x,y,x',y'\) & \(T\) ist der Träger der Baumstruktur.
  Jedes seiner Elemente \(x,y,x',y'\) steht für einen ganzen Baum.\\[3pt]
\(\ell\) & Der Blattkonstruktor ordnet einer Beschriftung \(a\)
  den einzelnen Blattbaum \(\ell(a)\) zu.\\[3pt]
\(k\) & Der Knotenkonstruktor verbindet zwei ganze Bäume \(x,y\)
  unter einer neuen Wurzel zu \(k(x,y)\): \(x\) ist der linke,
  \(y\) der rechte Teilbaum.\\[3pt]
\(T\times T\) & Die Menge der geordneten Paare ganzer Bäume;
  sie ist der Definitionsbereich von \(k\).\\[3pt]
\(U;\ \ell[A]\) & \(U\subseteq T\) bezeichnet eine beliebige
  Teilmenge des Baumträgers. \(\ell[A]\) ist die Menge aller
  durch \(\ell\) erzeugten Blattbäume.
\end{tabularx}
\end{center}

Im zuvor konstruierten Modell ist die Zuordnung konkret
\[
  T=\BracketTreeSet A,\qquad \ell=\lambda_A,\qquad k=\kappa_A.
\]
Die folgenden Bedingungen B0 bis B4 werden nun für die Daten
\((T,\ell,k)\) gefordert. Die jeweils genannten Herkunftssätze
weisen sie für dieses konkrete Modell nach. Der Ausdruck
\(\mathsf{BaumStr}_A(T,\ell,k)\) besagt, dass die gewählten Daten
alle fünf Bedingungen erfüllen.

\noindent\begin{minipage}{\linewidth}
\FormulaDefDeltaK[Freie volle planare Binärbaumstruktur]{%
  \mathsf{BaumStr}_A(T,\ell,k)
}{BinaryTreeStructureDef}{%
  \DeltaRow{Mengen}{A\dsep T\dsep U\dsep a\dsep b\dsep x\dsep y\dsep x'\dsep y'}
  \DeltaRow{Funktionen}{\ell\dsep k}
  \DeltaRow{\textbf{Strukturaxiome}}{}
  \DeltaRow{B0: Typisierung}{\ell\colon A\to T\land k\colon T\times T\to T}
  \DeltaRow{B1: Blatteindeutigkeit}{%
    \forall a,b\in A\;(\ell(a)=\ell(b)\rightarrow a=b)}
  \DeltaRow{B2: Knoteneindeutigkeit}{%
    \begin{aligned}[t]
      &\forall x,y,x',y'\in T\\[-2pt]
      &\quad(k(x,y)=k(x',y')\rightarrow x=x'\land y=y')
    \end{aligned}}
  \DeltaRow{B3: Trennung}{\forall a\in A\,\forall x,y\in T\;\ell(a)\ne k(x,y)}
  \DeltaRow{B4: Minimalität}{%
    \begin{aligned}[t]
      &\forall U\subseteq T\,\Bigl(\bigl(\ell[A]\subseteq U\land{}\\[-2pt]
      &\quad(\forall x,y\in U\;k(x,y)\in U)\bigr)\rightarrow U=T\Bigr)
    \end{aligned}}
  \DeltaRow{\textbf{Neues Symbol}}{}
  \DeltaPrem{Freie Binärbaumstruktur}{\mathsf{BaumStr}_A(T,\ell,k)}}
\end{minipage}\par

\Needspace{14\baselineskip}
\noindent\begin{minipage}{\linewidth}
\FormulaAxiomDeltaK[Typisierung einer Baumstruktur]{%
  \mathsf{BaumStr}_A(T,\ell,k)\vdash
    \ell\colon A\to T\land k\colon T\times T\to T
}{BinaryTreeStructureTypingAxiom}{%
  \DeltaRow{Mengen}{A\dsep T}\DeltaRow{Funktionen}{\ell\dsep k}}
\par\smallskip\noindent
B0 verlangt, dass beide Konstruktoren wieder Bäume in \(T\)
liefern. Für \(\lambda_A\) und \(\kappa_A\) auf dem konkreten
Baumträger wurde dies in
\FormulaRefAuto{TreeCodeConstructorTyping}[thm] bewiesen.
\end{minipage}\par

\Needspace{14\baselineskip}
\noindent\begin{minipage}{\linewidth}
\FormulaAxiomDeltaK[Injektivität des Blattkonstruktors]{%
  \mathsf{BaumStr}_A(T,\ell,k)\dsep a,b\in A\dsep\ell(a)=\ell(b)\vdash a=b
}{BinaryTreeStructureLeafAxiom}{%
  \DeltaRow{Mengen}{A\dsep T\dsep a\dsep b}\DeltaRow{Funktionen}{\ell\dsep k}}
\par\smallskip\noindent
B1 besagt: Gleiche Blattbäume haben dieselbe Beschriftung.
Für den konkreten Blattkonstruktor \(\lambda_A\) ist dies
\FormulaRefAuto{TreeCodeLeafInjectivity}[thm].
\end{minipage}\par

\Needspace{15\baselineskip}
\noindent\begin{minipage}{\linewidth}
\FormulaAxiomDeltaK[Injektivität des Knotenkonstruktors]{%
  \begin{aligned}[t]
    &\mathsf{BaumStr}_A(T,\ell,k)\dsep x,y,x',y'\in T\dsep{}\\[-2pt]
    &k(x,y)=k(x',y')\vdash x=x'\land y=y'
  \end{aligned}
}{BinaryTreeStructureNodeAxiom}{%
  \DeltaRow{Mengen}{A\dsep T\dsep x\dsep y\dsep x'\dsep y'}
  \DeltaRow{Funktionen}{\ell\dsep k}}
\par\smallskip\noindent
B2 besagt: Gleiche zusammengesetzte Bäume haben denselben linken
und denselben rechten Teilbaum. Für \(\kappa_A\) wurde dies in
\FormulaRefAuto{TreeCodeNodeInjectivity}[thm] bewiesen.
\end{minipage}\par

\Needspace{14\baselineskip}
\noindent\begin{minipage}{\linewidth}
\FormulaAxiomDeltaK[Trennung von Blättern und Knoten]{%
  \mathsf{BaumStr}_A(T,\ell,k)\dsep a\in A\dsep x,y\in T
    \vdash\ell(a)\ne k(x,y)
}{BinaryTreeStructureDisjointnessAxiom}{%
  \DeltaRow{Mengen}{A\dsep T\dsep a\dsep x\dsep y}
  \DeltaRow{Funktionen}{\ell\dsep k}}
\par\smallskip\noindent
B3 schließt aus, dass ein Blattbaum zugleich aus zwei Teilbäumen
zusammengesetzt ist. Die entsprechende Trennung von
\(\lambda_A\) und \(\kappa_A\) liefert
\FormulaRefAuto{TreeCodeConstructorSeparation}[thm].
\end{minipage}\par

\Needspace{16\baselineskip}
\noindent\begin{minipage}{\linewidth}
\FormulaAxiomDeltaK[Minimalität einer Baumstruktur]{%
  \begin{aligned}[t]
    &\mathsf{BaumStr}_A(T,\ell,k)\dsep U\subseteq T\dsep\ell[A]\subseteq U\dsep{}\\[-2pt]
    &\forall x,y\in U\;k(x,y)\in U\vdash U=T
  \end{aligned}
}{BinaryTreeStructureMinimalityAxiom}{%
  \DeltaRow{Mengen}{A\dsep T\dsep U\dsep x\dsep y}
  \DeltaRow{Funktionen}{\ell\dsep k}}
\par\smallskip\noindent
B4 besagt: Enthält \(U\) alle Blattbäume \(\ell[A]\) und mit
\(x,y\) jeweils auch \(k(x,y)\), so enthält \(U\) bereits alle
Bäume von \(T\). Diese Minimalität des konkreten Baumträgers
ist \FormulaRefAuto{TreeCodeMinimality}[thm].
\end{minipage}\par

Die zuvor bewiesenen Eigenschaften werden nun in einem Strukturbegriff
gebündelt. Die registrierten Axiome sind Zugriffsregeln auf die
Bedingungen B0 bis B4; sie fordern keine neuen Mengenexistenzaxiome.

\noindent\begin{minipage}{\linewidth}
\FormulaThmDeltaK[Existenz einer freien Klammerungsbaumstruktur]{%
  \mathsf{BaumStr}_A(\BracketTreeSet A,\lambda_A,\kappa_A)
}{BinaryTreeStructureConcreteModel}{%
  \DeltaRow{Mengen}{A\dsep U\dsep S\dsep R\dsep S'\dsep R'\dsep a\dsep b}}
\end{minipage}\par
\begin{tabproofwide}
  \proofstepwidestar[]{\begin{aligned}[t]&\lambda_A\colon A\to\BracketTreeSet A\land{}\\&\kappa_A\colon\BracketTreeSet A\times\BracketTreeSet A\to\BracketTreeSet A\end{aligned}}{%
    \FormulaRefAuto{TreeCodeConstructorTyping}[thm]}
  \proofstepwidestar[2]{a\in A}{%
    \rA}
  \proofstepwidestar[3]{b\in A}{%
    \rA}
  \proofstepwidestar[4]{\lambda_A(a)=\lambda_A(b)}{%
    \rA}
  \proofstepwidestar[2,3,4]{a=b}{%
    \FormulaRefAuto{TreeCodeLeafInjectivity}[thm]{\rAI{2,3},4}}
  \proofstepwidestar[]{\forall a,b\in A\;(\lambda_A(a)=\lambda_A(b)\rightarrow a=b)}{%
    \rUI{\rRI{2,\rUI{\rRI{3,\rRI{4,5}}}}}}
  \proofstepwidestar[7]{S\in\BracketTreeSet A}{%
    \rA}
  \proofstepwidestar[8]{R\in\BracketTreeSet A}{%
    \rA}
  \proofstepwidestar[9]{S'\in\BracketTreeSet A}{%
    \rA}
  \proofstepwidestar[10]{R'\in\BracketTreeSet A}{%
    \rA}
  \proofstepwidestar[11]{\kappa_A(S,R)=\kappa_A(S',R')}{%
    \rA}
  \proofstepwidestar[7,8,9,10,11]{S=S'\land R=R'}{%
    \FormulaRefAuto{TreeCodeNodeInjectivity}[thm]{\rAI{7,\rAI{8,\rAI{9,10}}},11}}
  \proofstepwidestar[]{\begin{aligned}[t]&\forall S,R,S',R'\in\BracketTreeSet A\\&(\kappa_A(S,R)=\kappa_A(S',R')\rightarrow S=S'\land R=R')\end{aligned}}{%
    \rUI{\rRI{7,\rUI{\rRI{8,\rUI{\rRI{9,\rUI{\rRI{10,\rRI{11,12}}}}}}}}}}
  \proofstepwidestar[14]{a\in A}{%
    \rA}
  \proofstepwidestar[15]{S\in\BracketTreeSet A}{%
    \rA}
  \proofstepwidestar[16]{R\in\BracketTreeSet A}{%
    \rA}
  \proofstepwidestar[14,15,16]{\lambda_A(a)\ne\kappa_A(S,R)}{%
    \FormulaRefAuto{TreeCodeConstructorSeparation}[thm]{14,\rAI{15,16}}}
  \proofstepwidestar[]{\forall a\in A\,\forall S,R\in\BracketTreeSet A\;\lambda_A(a)\ne\kappa_A(S,R)}{%
    \rUI{\rRI{14,\rUI{\rRI{15,\rUI{\rRI{16,17}}}}}}}
  \proofstepwidestar[19]{U\subseteq\BracketTreeSet A}{%
    \rA}
  \proofstepwidestar[20]{\lambda_A[A]\subseteq U\land(\forall S,R\in U\;\kappa_A(S,R)\in U)}{%
    \rA}
  \proofstepwidestar[19,20]{U=\BracketTreeSet A}{%
    \FormulaRefAuto{TreeCodeMinimality}[thm]{19,\rAEa{20},\rAEb{20}}}
  \proofstepwidestar[]{\begin{aligned}[t]&\forall U\subseteq\BracketTreeSet A\;\Bigl((\lambda_A[A]\subseteq U\land{}\\&\quad(\forall S,R\in U\;\kappa_A(S,R)\in U))\rightarrow U=\BracketTreeSet A\Bigr)\end{aligned}}{%
    \rUI{\rRI{19,\rRI{20,21}}}}
  \proofstepwidestar[]{\mathsf{BaumStr}_A(\BracketTreeSet A,\lambda_A,\kappa_A)}{%
    \FormulaRefAuto{BinaryTreeStructureDef}[def]{1,6,13,18,22}}
\end{tabproofwide}

\subsection{Erste Folgerungen aus den Baumaxiomen}

Von hier an werden Induktion und Zerlegung aus den Strukturaxiomen
hergeleitet. Die einzelnen Codepositionen und Baumstufen gehen in diese
Beweise nicht mehr ein. Die gewonnenen Sätze gelten deshalb für jede
Menge mit Konstruktoren, welche B0 bis B4 erfüllt.

\noindent\begin{minipage}{\linewidth}
\FormulaThmDeltaK[Blätter liegen im Strukturträger]{%
\mathsf{BaumStr}_A(T,\ell,k)\dsep a\in A\vdash\ell(a)\in T
}{BinaryTreeStructureLeafTyping}{%
\DeltaRow{Mengen}{A\dsep T\dsep a}\DeltaRow{Funktionen}{\ell\dsep k}}
\end{minipage}\par
\begin{tabproofwide}
  \proofstepwidestar[1]{\mathsf{BaumStr}_A(T,\ell,k)}{%
    \rA}
  \proofstepwidestar[2]{a\in A}{%
    \rA}
  \proofstepwidestar[1]{\ell\colon A\to T}{%
    \rAEa{\FormulaRefAuto{BinaryTreeStructureTypingAxiom}[ax]{1}}}
  \proofstepwidestar[1,2]{\ell(a)\in T}{%
    \FormulaRefAuto{F\colon A\to B\dsep x\in A\vdash F(x)\in B}[thm]{3,2}}
\end{tabproofwide}

\noindent\begin{minipage}{\linewidth}
\FormulaThmDeltaK[Knoten liegen im Strukturträger]{%
\mathsf{BaumStr}_A(T,\ell,k)\dsep x,y\in T\vdash k(x,y)\in T
}{BinaryTreeStructureNodeTyping}{%
\DeltaRow{Mengen}{A\dsep T\dsep x\dsep y}\DeltaRow{Funktionen}{\ell\dsep k}}
\end{minipage}\par
\begin{tabproofwide}
  \proofstepwidestar[1]{\mathsf{BaumStr}_A(T,\ell,k)}{%
    \rA}
  \proofstepwidestar[2]{x,y\in T}{%
    \rA}
  \proofstepwidestar[1]{k\colon T\times T\to T}{%
    \rAEb{\FormulaRefAuto{BinaryTreeStructureTypingAxiom}[ax]{1}}}
  \proofstepwidestar[2]{(x,y)\in T\times T}{%
    \FormulaRefAuto{a\in A\dsep b\in B\vdash(a,b)\in A\times B}[thm]{\rAEa{2},\rAEb{2}}}
  \proofstepwidestar[1,2]{k(x,y)\in T}{%
    \FormulaRefAuto{F\colon A\to B\dsep x\in A\vdash F(x)\in B}[thm]{3,4}}
\end{tabproofwide}

\noindent\begin{minipage}{\linewidth}
\FormulaThmDeltaK[Eindeutigkeit des geordneten Kinderpaares]{%
\mathsf{BaumStr}_A(T,\ell,k)\dsep p,q\in T\times T\dsep k(p)=k(q)\vdash p=q
}{BinaryTreeStructurePairInjective}{%
\DeltaRow{Mengen}{A\dsep T\dsep p\dsep q\dsep x\dsep y\dsep x'\dsep y'}\DeltaRow{Funktionen}{\ell\dsep k}}
\end{minipage}\par
\begin{tabproofwide}
  \proofstepwidestar[1]{\mathsf{BaumStr}_A(T,\ell,k)}{%
    \rA}
  \proofstepwidestar[2]{p,q\in T\times T}{%
    \rA}
  \proofstepwidestar[3]{k(p)=k(q)}{%
    \rA}
  \proofstepwidestar[2]{\pi_1(p)\in T}{%
    \FormulaRefAuto{z\in A\times B\vdash\pi_1(z)\in A}[thm]{\rAEa{2}}}
  \proofstepwidestar[2]{\pi_2(p)\in T}{%
    \FormulaRefAuto{z\in A\times B\vdash\pi_2(z)\in B}[thm]{\rAEa{2}}}
  \proofstepwidestar[2]{\pi_1(q)\in T}{%
    \FormulaRefAuto{z\in A\times B\vdash\pi_1(z)\in A}[thm]{\rAEb{2}}}
  \proofstepwidestar[2]{\pi_2(q)\in T}{%
    \FormulaRefAuto{z\in A\times B\vdash\pi_2(z)\in B}[thm]{\rAEb{2}}}
  \proofstepwidestar[2]{p=(\pi_1(p),\pi_2(p))}{%
    \FormulaRefAuto{z\in A\times B\vdash z=\bigl(\pi_1(z),\pi_2(z)\bigr)}[thm]{\rAEa{2}}}
  \proofstepwidestar[2]{q=(\pi_1(q),\pi_2(q))}{%
    \FormulaRefAuto{z\in A\times B\vdash z=\bigl(\pi_1(z),\pi_2(z)\bigr)}[thm]{\rAEb{2}}}
  \proofstepwidestar[2,3]{k(\pi_1(p),\pi_2(p))=k(\pi_1(q),\pi_2(q))}{%
    \rIE{8,\rIE{9,3}}}
  \proofstepwidestar[1,2,3]{\pi_1(p)=\pi_1(q)\land\pi_2(p)=\pi_2(q)}{%
    \FormulaRefAuto{BinaryTreeStructureNodeAxiom}[ax]{1,\rAI{4,\rAI{5,\rAI{6,7}}},10}}
  \proofstepwidestar[1,2,3]{(\pi_1(p),\pi_2(p))=(\pi_1(q),\pi_2(q))}{%
    \FormulaRefAuto{(a,b)=(c,d)\eqvdash a=c\land b=d}[thm]{11}}
  \proofstepwidestar[1,2,3]{p=q}{%
    \rIE{\FormulaRefAuto{a=b\vdash b=a}[thm]{8},\rIE{\FormulaRefAuto{a=b\vdash b=a}[thm]{9},12}}}
\end{tabproofwide}

\noindent\begin{minipage}{\linewidth}
\FormulaThmDeltaK[Induktion in einer Baumstruktur]{%
  \begin{aligned}[t]
    &\mathsf{BaumStr}_A(T,\ell,k)\dsep\forall a\in A\;P(\ell(a))\dsep{}\\[-2pt]
    &\forall x,y\in T\;(P(x)\land P(y)\rightarrow P(k(x,y)))
      \vdash\forall t\in T\;P(t)
  \end{aligned}
}{BinaryTreeStructureInduction}{%
  \DeltaRow{Mengen}{A\dsep T\dsep U\dsep a\dsep x\dsep y\dsep t\dsep z}
  \DeltaRow{Funktionen}{\ell\dsep k}\DeltaRow{Prädikat}{P}
  \DeltaRow{Abkürzung}{U\coloneqq\{t\in T\mid P(t)\}}}
\end{minipage}\par
\begin{tabproofwide}
  \proofstepwidestar[1]{\mathsf{BaumStr}_A(T,\ell,k)}{%
    \rA}
  \proofstepwidestar[2]{\forall a\in A\;P(\ell(a))}{%
    \rA}
  \proofstepwidestar[3]{\forall x,y\in T\;(P(x)\land P(y)\rightarrow P(k(x,y)))}{%
    \rA}
  \proofstepwidestar[]{U\subseteq T}{%
    \FormulaRefAuto{\{x\in A\mid P(x)\}\subseteq A}[thm]}
  \proofstepwidestar[1]{\ell\colon A\to T}{%
    \rAEa{\FormulaRefAuto{BinaryTreeStructureTypingAxiom}[ax]{1}}}
  \proofstepwidestar[6]{z\in\ell[A]}{%
    \rA}
  \proofstepwidestar[1,6]{\exists a\in A\;z=\ell(a)}{%
    \FormulaRefAuto{F\colon A\to B\dsep y\in F[A]\vdash\exists x\in A\;y=F(x)}[thm]{5,6}}
  \proofstepwidestar[8]{a\in A\land z=\ell(a)}{%
    \rA}
  \proofstepwidestar[1,8]{\ell(a)\in T}{%
    \FormulaRefAuto{BinaryTreeStructureLeafTyping}[thm]{1,\rAEa{8}}}
  \proofstepwidestar[2,8]{P(\ell(a))}{%
    \rRE{\rUE{2},\rAEa{8}}}
  \proofstepwidestar[1,2,8]{\ell(a)\in U}{%
    \FormulaRefAuto{x\in A\dsep P(x)\vdash x\in\{x\in A\mid P(x)\}}[thm]{9,10}}
  \proofstepwidestar[1,2,8]{z\in U}{%
    \rIE{\FormulaRefAuto{a=b\vdash b=a}[thm]{\rAEb{8}},11}}
  \proofstepwidestar[1,2,6]{z\in U}{%
    \rEE{7,8,12}}
  \proofstepwidestar[1,2]{\ell[A]\subseteq U}{%
    \FormulaRefAuto{SubsetIntroductionRule}[thm]{[6]\vdots 13}}
  \proofstepwidestar[15]{x\in U}{%
    \rA}
  \proofstepwidestar[16]{y\in U}{%
    \rA}
  \proofstepwidestar[15,16]{x\in T\land P(x)}{%
    \FormulaRefAuto{x\in\{u\in A\mid P(u)\}\eqvdash x\in A\land P(x)}[thm]{15}}
  \proofstepwidestar[15,16]{y\in T\land P(y)}{%
    \FormulaRefAuto{x\in\{u\in A\mid P(u)\}\eqvdash x\in A\land P(x)}[thm]{16}}
  \proofstepwidestar[1,15,16]{k(x,y)\in T}{%
    \FormulaRefAuto{BinaryTreeStructureNodeTyping}[thm]{1,\rAI{\rAEa{17},\rAEa{18}}}}
  \proofstepwidestar[3,15,16]{P(k(x,y))}{%
    \rRE{\rRE{\rUE{\rRE{\rUE{3},\rAEa{17}}},\rAEa{18}},\rAI{\rAEb{17},\rAEb{18}}}}
  \proofstepwidestar[1,3,15,16]{k(x,y)\in U}{%
    \FormulaRefAuto{x\in A\dsep P(x)\vdash x\in\{x\in A\mid P(x)\}}[thm]{19,20}}
  \proofstepwidestar[1,3]{\forall x,y\in U\;k(x,y)\in U}{%
    \rUI{\rRI{15,\rUI{\rRI{16,21}}}}}
  \proofstepwidestar[1,2,3]{U=T}{%
    \FormulaRefAuto{BinaryTreeStructureMinimalityAxiom}[ax]{1,4,14,22}}
  \proofstepwidestar[24]{t\in T}{%
    \rA}
  \proofstepwidestar[1,2,3,24]{t\in U}{%
    \rIE{\FormulaRefAuto{a=b\vdash b=a}[thm]{23},24}}
  \proofstepwidestar[1,2,3,24]{P(t)}{%
    \rAEb{\FormulaRefAuto{x\in\{u\in A\mid P(u)\}\eqvdash x\in A\land P(x)}[thm]{25}}}
  \proofstepwidestar[1,2,3]{\forall t\in T\;P(t)}{%
    \rUI{\rRI{24,26}}}
\end{tabproofwide}

Auch hier folgt die Minimalität umgekehrt durch Einsetzen von
\(P(t)\coloneqq t\in U\). Insbesondere schließt das Induktionsaxiom
zusätzliche Elemente ohne endliche Konstruktion aus.

\noindent\begin{minipage}{\linewidth}
\FormulaThmDeltaK[Eindeutige Zerlegung in einer Baumstruktur]{%
  \begin{aligned}[t]
    &\mathsf{BaumStr}_A(T,\ell,k)\dsep t\in T\vdash{}\\[-2pt]
    &\bigl(\exists!a\in A\;t=\ell(a)\bigr)
      \lor\bigl(\exists!p\in T\times T\;t=k(p)\bigr)
  \end{aligned}
}{BinaryTreeStructureDecomposition}{%
  \DeltaRow{Mengen}{A\dsep T\dsep t\dsep a\dsep b\dsep c\dsep p\dsep q\dsep x\dsep y}
  \DeltaRow{Funktionen}{\ell\dsep k}
  \DeltaRow{Prädikatsabkürzung}{\begin{aligned}[t]
    Q(t)\coloneqq{}&(\exists!a\in A\;t=\ell(a))\lor{}\\[-2pt]
    &(\exists!p\in T\times T\;t=k(p))
  \end{aligned}}}
\end{minipage}\par
\begin{tabproofwide}
  \proofstepwidestar[1]{\mathsf{BaumStr}_A(T,\ell,k)}{%
    \rA}
  \proofstepwidestar[2]{t\in T}{%
    \rA}
  \proofstepwidestar[3]{a\in A}{%
    \rA}
  \proofstepwidestar[3]{\exists b\in A\;\ell(a)=\ell(b)}{%
    \rEI{\rAI{3,\rII}}}
  \proofstepwidestar[5]{b\in A\land\ell(a)=\ell(b)}{%
    \rA}
  \proofstepwidestar[6]{c\in A\land\ell(a)=\ell(c)}{%
    \rA}
  \proofstepwidestar[5,6]{\ell(b)=\ell(c)}{%
    \rIE{\rAEb{5},\rAEb{6}}}
  \proofstepwidestar[1,5,6]{b=c}{%
    \FormulaRefAuto{BinaryTreeStructureLeafAxiom}[ax]{1,\rAEa{5},\rAEa{6},7}}
  \proofstepwidestar[1,3]{\exists!b\in A\;\ell(a)=\ell(b)}{%
    \UEI{4,5,6,8}}
  \proofstepwidestar[1,3]{Q(\ell(a))}{%
    \rOIa{9}}
  \proofstepwidestar[1]{\forall a\in A\;Q(\ell(a))}{%
    \rUI{\rRI{3,10}}}
  \proofstepwidestar[12]{x\in T}{%
    \rA}
  \proofstepwidestar[13]{y\in T}{%
    \rA}
  \proofstepwidestar[12,13]{(x,y)\in T\times T}{%
    \FormulaRefAuto{a\in A\dsep b\in B\vdash(a,b)\in A\times B}[thm]{12,13}}
  \proofstepwidestar[12,13]{\exists p\in T\times T\;k(x,y)=k(p)}{%
    \rEI{\rAI{14,\rII}}}
  \proofstepwidestar[16]{p\in T\times T\land k(x,y)=k(p)}{%
    \rA}
  \proofstepwidestar[17]{q\in T\times T\land k(x,y)=k(q)}{%
    \rA}
  \proofstepwidestar[16,17]{k(p)=k(q)}{%
    \rIE{\rAEb{16},\rAEb{17}}}
  \proofstepwidestar[1,16,17]{p=q}{%
    \FormulaRefAuto{BinaryTreeStructurePairInjective}[thm]{1,\rAI{\rAEa{16},\rAEa{17}},18}}
  \proofstepwidestar[1,12,13]{\exists!p\in T\times T\;k(x,y)=k(p)}{%
    \UEI{15,16,17,19}}
  \proofstepwidestar[1,12,13]{Q(k(x,y))}{%
    \rOIb{20}}
  \proofstepwidestar[22]{Q(x)\land Q(y)}{%
    \rA}
  \proofstepwidestar[1,12,13]{Q(x)\land Q(y)\rightarrow Q(k(x,y))}{%
    \rRI{22,21}}
  \proofstepwidestar[1]{\forall x,y\in T\;(Q(x)\land Q(y)\rightarrow Q(k(x,y)))}{%
    \rUI{\rRI{12,\rUI{\rRI{13,23}}}}}
  \proofstepwidestar[1]{\forall t\in T\;Q(t)}{%
    \FormulaRefAuto{BinaryTreeStructureInduction}[thm]{1,11,24}}
  \proofstepwidestar[1,2]{(\exists!a\in A\;t=\ell(a))\lor(\exists!p\in T\times T\;t=k(p))}{%
    \rRE{\rUE{25},2}}
\end{tabproofwide}

Die beiden Fälle sind nach
\FormulaRefAuto{BinaryTreeStructureDisjointnessAxiom}[ax] disjunkt.

